% ============================================================
%  Introducing Python, chapter 4: Strings
% ============================================================
\chapter{Strings}

Most programs spend more time on text than on math, and AI work is text work: prompts,
model replies, file names, JSON. This chapter covers how to create strings, pick them apart
with indexes and slices, search and clean them, and build new ones with f-strings. Everything here
applies to prompts you will format and replies you will clean.

% ------------------------------------------------------------
\section{Strings Are Sequences}

A \code{str} is an ordered \textbf{sequence} of characters. A character is the smallest unit of
writing: a letter, digit, symbol, space or newline. Python has no separate character type.
A single character is just a string of length 1, and a string may be empty.

Strings are \textbf{immutable}: no operation edits a string in place. Every operation
that seems to change one returns a \emph{new} string. This is the most important fact in the chapter.

Every sequence (strings, and later lists and tuples) supports the same core operations:

\begin{center}
\begin{tabularx}{\linewidth}{@{}l l l X@{}}
\toprule
\textbf{Operation} & \textbf{Example} & \textbf{Returns} & \textbf{Meaning} \\
\midrule
membership & \code{"fun" in s} & \code{bool} & is it present? \\
position of a value & \code{s.index("x")} & \code{int} & offset of first match \\
one element & \code{s[0]} & \code{str} (length 1) & item at an offset \\
slice & \code{s[2:5]} & \code{str} & a range of items \\
length & \code{len(s)} & \code{int} & number of items \\
smallest, largest & \code{min(s)}, \code{max(s)} & \code{str} (length 1) & by character code \\
\bottomrule
\end{tabularx}
\end{center}

\begin{pycode}
>>> prompt = "Summarize this text"
>>> "text" in prompt
True
>>> prompt.index("this")
10
>>> prompt[0]
'S'
>>> prompt[10:14]
'this'
>>> len(prompt)
19
>>> min("model"), max("model")
('d', 'o')
\end{pycode}

% ------------------------------------------------------------
\section{Create with Quotes}

Wrap characters in matching single or double quotes. Python treats both the same.
The interpreter echoes strings with single quotes.

\begin{pycode}
>>> 'single' == "double"
False
>>> 'tokens' == "tokens"
True
>>> "It's fine"
"It's fine"
>>> 'She said "hi"'
'She said "hi"'
\end{pycode}

Having two quote styles lets you put one kind of quote inside the other. Triple quotes
(\code{'''} or \code{"""}) create \emph{multi-line} strings and keep the line breaks.
That is the usual way to write a long prompt.

\begin{pycode}
>>> text = """line one
... line two"""
>>> text
'line one\nline two'
>>> print(text)
line one
line two
\end{pycode}

\paragraph{Echo versus print.} The interpreter echo shows the string's representation, with quotes and
escape codes like \code{\textbackslash n}. \code{print()} shows the text itself, separates its arguments
with spaces and ends with a newline.

\begin{pycode}
>>> print("a", "b", 3)
a b 3
>>> len(""), bool("")
(0, False)
\end{pycode}

\paragraph{Prefixes.} A letter before the opening quote changes how the string is read.

\begin{center}
\begin{tabularx}{\linewidth}{@{}l l X@{}}
\toprule
\textbf{Prefix} & \textbf{Kind} & \textbf{Use} \\
\midrule
none & plain text (\code{str}) & most strings \\
\code{f} & f-string & insert values: \code{f"hi \{name\}"} \\
\code{r} & raw string & keep backslashes literally (paths, regular expressions) \\
\code{fr} & raw f-string & both \\
\code{b} & \code{bytes}, not \code{str} & binary data (Chapter 5) \\
\code{u} & same as plain & historical; you can ignore it \\
\bottomrule
\end{tabularx}
\end{center}

\section{Create with str()}

\code{str(x)} converts any object to its text form. \code{print()} and f-strings do this for you.
\code{+} does \emph{not}: mixing \code{str} and \code{int} with \code{+} is a \code{TypeError}.

\begin{pycode}
>>> str(42), str(3.5), str(True), str(None)
('42', '3.5', 'True', 'None')
>>> str([1, 2])
'[1, 2]'
>>> "tokens: " + 42
Traceback (most recent call last):
  ...
TypeError: can only concatenate str (not "int") to str
>>> "tokens: " + str(42)
'tokens: 42'
\end{pycode}

\section{Escape with \textbackslash}

A backslash gives the next character a special meaning (an \textbf{escape sequence}), or
removes its special meaning.

\begin{center}
\begin{tabular}{@{}l l@{}}
\toprule
\textbf{Write} & \textbf{You get} \\
\midrule
\code{\textbackslash n} & newline \\
\code{\textbackslash t} & tab \\
\code{\textbackslash \textbackslash} & one backslash \\
\code{\textbackslash '} and \code{\textbackslash "} & a quote character \\
\bottomrule
\end{tabular}
\end{center}

\begin{pycode}
>>> "a\tb"
'a\tb'
>>> print("line1\nline2")
line1
line2
>>> print('C:\\logs\\new')
C:\logs\new
\end{pycode}

A \textbf{raw string} (\code{r"..."}) turns escapes off. It does not affect real line breaks inside triple quotes.

\begin{pycode}
>>> print(r"C:\logs\new")
C:\logs\new
>>> len(r"\n"), len("\n")
(2, 1)
\end{pycode}

\begin{trap}
On Windows a path like \code{"C:\textbackslash new"} contains a newline, because \code{\textbackslash n} is
an escape. Use a raw string, double the backslashes, or use forward slashes.
\begin{pycode}
>>> print("C:\new")
C:
ew
\end{pycode}
\end{trap}

\section{Combine with + and Duplicate with *}

\code{+} joins two strings; \code{*} repeats a string by an \code{int}. Python adds no spaces.
Two string \emph{literals} next to each other are joined automatically, which is handy for splitting
a long string over several lines inside parentheses.

\begin{pycode}
>>> "ab" + "cd"
'abcd'
>>> "ab" "cd"
'abcd'
>>> "-" * 10
'----------'
>>> "ha" * 3 + "!"
'hahaha!'
>>> "5" + "5", "5" * 3
('55', '555')
>>> prompt = ("You are a helpful assistant. "
...           "Answer in one sentence.")
>>> prompt
'You are a helpful assistant. Answer in one sentence.'
\end{pycode}

\code{*} binds tighter than \code{+}, so \code{"ha" * 3 + "!"} repeats first.

% ------------------------------------------------------------
\section{Get a Character by [offset]}

Offsets start at 0 on the left. Negative offsets count from the right, so \code{-1} is the last character.

%MD
%```text
%  word =   t    o    k    e    n    i    z    e
%  offset:  0    1    2    3    4    5    6    7
%  offset: -8   -7   -6   -5   -4   -3   -2   -1
%```
%ENDMD
\begin{center}
\begin{tikzpicture}[font=\ttfamily\small]
  \foreach \q [count=\i from 0] in {t,o,k,e,n,i,z,e} {
    \pgfmathtruncatemacro{\neg}{\i-8}
    \node[draw=cblue,fill=cblue!8,minimum width=9mm,minimum height=9mm] at (1.0*\i,0) {\q};
    \node[font=\sffamily\scriptsize,text=cgreen] at (1.0*\i,0.7) {\i};
    \node[font=\sffamily\scriptsize,text=cred] at (1.0*\i,-0.7) {\neg};
  }
  \node[anchor=east,font=\sffamily\scriptsize,text=cgreen] at (-0.7,0.7) {offset};
  \node[anchor=east,font=\sffamily\scriptsize,text=cred] at (-0.7,-0.7) {from the end};
\end{tikzpicture}
\end{center}

\begin{pycode}
>>> word = "tokenize"
>>> word[0], word[3], word[-1], word[-2]
('t', 'e', 'e', 'z')
>>> word[8]
Traceback (most recent call last):
  ...
IndexError: string index out of range
\end{pycode}

An offset past the end raises \code{IndexError}. You cannot assign to an offset, because strings are immutable.
Build a new string instead:

\begin{pycode}
>>> word = "tokenize"
>>> word[0] = "T"
Traceback (most recent call last):
  ...
TypeError: 'str' object does not support item assignment
>>> "T" + word[1:]
'Tokenize'
>>> word.replace("t", "T", 1)
'Tokenize'
>>> word
'tokenize'
\end{pycode}

\begin{typecheck}[: what indexing returns]
\code{s[i]} returns a \code{str} of length 1, never a separate char type. \code{s[i]} needs an
\code{int}; a \code{str} offset such as \code{s["a"]} is a \code{TypeError}. Compare with a dict, where the subscript
is a key, not a position.
\end{typecheck}

\section{Get a Substring with a Slice}

\code{s[start:stop:step]} returns a new string. Think of the offsets as numbers on the
\emph{lines between characters}. \code{start} is included; \code{stop} is excluded. Omitted parts
default to the beginning, the end, and step 1.

%MD
%```text
%   cut lines:  0   1   2   3   4   5   6   7   8   9   10
%              +---+---+---+---+---+---+---+---+---+---+
%   string:    | 0 | 1 | 2 | 3 | 4 | 5 | 6 | 7 | 8 | 9 |
%              +---+---+---+---+---+---+---+---+---+---+
%                      \___________/
%                        s[2:5] = "234"
%```
%ENDMD
\begin{center}
\begin{tikzpicture}[font=\ttfamily\small]
  \foreach \i in {0,...,9} {
    \ifnum\i>1 \ifnum\i<5 \def\fillc{cgold!45}\else\def\fillc{cblue!8}\fi\else\def\fillc{cblue!8}\fi
    \node[draw=cblue,fill=\fillc,minimum width=9mm,minimum height=9mm] at (0.9*\i+0.45,0) {\i};
  }
  \foreach \k in {0,...,10} {
    \draw[cgray] (0.9*\k,-0.5) -- (0.9*\k,-0.62);
    \node[font=\sffamily\scriptsize,text=cgreen] at (0.9*\k,-0.85) {\k};
  }
  \draw[decorate,decoration={brace,amplitude=4pt},cgold!80!black] (1.8,0.55) -- (4.5,0.55);
  \node[font=\sffamily\scriptsize,text=cgold!70!black] at (3.15,1.1) {\texttt{s[2:5]} gives \texttt{"234"}};
\end{tikzpicture}
\end{center}

\begin{pycode}
>>> d = "0123456789"
>>> d[2:5], d[:3], d[7:], d[:]
('234', '012', '789', '0123456789')
>>> d[-3:], d[:-3], d[-5:-2]
('789', '0123456', '567')
>>> d[::2], d[1::2], d[2:9:3]
('02468', '13579', '258')
>>> d[::-1], d[8:2:-2]
('9876543210', '864')
\end{pycode}

\begin{center}
\begin{tabularx}{\linewidth}{@{}l l X@{}}
\toprule
\textbf{Slice} & \textbf{Result for \code{d}} & \textbf{Reading it} \\
\midrule
\code{d[2:5]} & \code{'234'} & from cut 2 up to (not including) cut 5 \\
\code{d[:3]} & \code{'012'} & first three characters \\
\code{d[7:]} & \code{'789'} & from cut 7 to the end \\
\code{d[-3:]} & \code{'789'} & last three characters \\
\code{d[:-3]} & \code{'0123456'} & everything except the last three \\
\code{d[::2]} & \code{'02468'} & every second character \\
\code{d[::-1]} & \code{'9876543210'} & reversed \\
\bottomrule
\end{tabularx}
\end{center}

Slices are forgiving. Out-of-range offsets are clamped instead of raising an error. A slice
that selects nothing gives the empty string. Indexing (\code{d[10]}) is strict; slicing is not.

\begin{pycode}
>>> d = "0123456789"
>>> d[-50:3], d[7:99], d[20:30], d[5:2]
('012', '789', '', '')
>>> d[10]
Traceback (most recent call last):
  ...
IndexError: string index out of range
\end{pycode}

\begin{boss}[: slices and cut lines]
Most slice mistakes come from counting characters instead of cut lines. Rules that always hold:
\begin{itemize}[itemsep=1pt]
\item \code{s[a:b]} has length \code{b - a} when both are in range.
\item \code{s[:n] + s[n:] == s} for any \code{n}. The two halves meet at one cut.
\item With a negative step, start and stop swap roles: you walk backward from \code{start} down to, but not including, \code{stop}.
\item Prefer simple forms: \code{s[:n]}, \code{s[-n:]}, \code{s[::-1]}.
\end{itemize}
%MD
%```mermaid
%flowchart LR
%    S["s[start:stop:step]"] --> A["start: first kept (included)"]
%    S --> B["stop: first NOT kept (excluded)"]
%    S --> C["step: stride, negative walks backward"]
%```
%ENDMD
\begin{center}
\begin{tikzpicture}[font=\sffamily\small,>=Stealth,
  b/.style={draw=cblue,fill=cblue!8,rounded corners=2pt,inner sep=4pt,align=center}]
  \node[font=\ttfamily\Large] (s) at (5,1.4) {s[start:stop:step]};
  \node[b] (a) at (0.8,-0.4) {start\\[-2pt]{\scriptsize first one kept}};
  \node[b] (b2) at (5,-0.4) {stop\\[-2pt]{\scriptsize first one NOT kept}};
  \node[b] (c) at (9.4,-0.4) {step\\[-2pt]{\scriptsize negative: walk backward}};
  \draw[->,thick,cgray] (s) -- (a);
  \draw[->,thick,cgray] (s) -- (b2);
  \draw[->,thick,cgray] (s) -- (c);
\end{tikzpicture}
\end{center}
\begin{pycode}
>>> d = "0123456789"
>>> d[:4] + d[4:] == d
True
>>> len(d[3:8])
5
>>> d[:3] + "!" + d[-2:]
'012!89'
\end{pycode}
\end{boss}

\section{Get Length with len()}

\code{len(s)} returns an \code{int}: the number of characters, including spaces and newlines.
The last valid offset is \code{len(s) - 1}. (Token counts are not character counts; models read
tokens, which are chunks of characters.)

\begin{pycode}
>>> len("hello world"), len("a\nb"), len("")
(11, 3, 0)
\end{pycode}

% ------------------------------------------------------------
\section{Methods: Functions That Belong to a String}

Many string operations are \textbf{methods}, called with a dot: \code{s.method(args)}. The dot
(attribute access) is different from the brackets used for indexing. A method \emph{never}
changes \code{s}; it returns a result.

\begin{typecheck}[: what the common methods return]
\begin{center}
\begin{tabularx}{\linewidth}{@{}l l X@{}}
\toprule
\textbf{Methods} & \textbf{Returns} & \textbf{Examples} \\
\midrule
transforming & new \code{str} & \code{upper lower strip replace title center} \\
splitting & \code{list[str]} & \code{split splitlines} \\
joining & new \code{str} & \code{sep.join(list)} \\
counting, finding & \code{int} & \code{find index count} \\
testing & \code{bool} & \code{startswith endswith isdigit isalpha isalnum} \\
\bottomrule
\end{tabularx}
\end{center}
\end{typecheck}

\begin{trap}
Forgetting to keep the result is the most common beginner bug with strings. Calling
\code{name.upper()} alone computes a new string and throws it away. \code{name} is untouched.
%MD
%```mermaid
%flowchart LR
%    N1["name -> 'ada'"] -->|"name.upper()"| X["'ADA' (new object, no name)"]
%    X -->|"name = ..."| N2["name -> 'ADA'"]
%```
%ENDMD
\begin{center}
\begin{tikzpicture}[font=\ttfamily\small,>=Stealth,
  ob/.style={draw=cblue,fill=cblue!8,inner sep=4pt}]
  \node[draw=cgold,fill=cgold!25,inner sep=3pt] (n) at (0,0) {name};
  \node[ob] (a) at (3,0) {"ada"};
  \draw[->,thick] (n) -- (a);
  \node[ob,draw=cgreen,fill=cgreen!10] (b) at (7.2,0) {"ADA"};
  \draw[->,thick,cgray] (a) -- node[above,font=\scriptsize]{name.upper()} node[below,font=\sffamily\scriptsize]{new object} (b);
  \node[draw=cgold,fill=cgold!25,inner sep=3pt] (n2) at (10.5,0) {name};
  \draw[->,thick,dashed,cgray] (n2) -- node[above,font=\sffamily\scriptsize]{only if} node[below,font=\scriptsize]{name = ...} (b);
\end{tikzpicture}
\end{center}
\begin{pycode}
>>> name = "ada"
>>> name.upper()
'ADA'
>>> name
'ada'
>>> name = name.upper()
>>> name
'ADA'
\end{pycode}
\end{trap}

\subsection{Split with split()}

\code{s.split(sep)} cuts a string at each separator and returns a \code{list} of strings. With no
argument it splits on any run of whitespace and drops empty pieces. A second argument limits the
number of cuts.

\begin{pycode}
>>> line = "role:user,content:hello,tokens:5"
>>> line.split(",")
['role:user', 'content:hello', 'tokens:5']
>>> line.split(",", 1)
['role:user', 'content:hello,tokens:5']
>>> "a  b\tc\nd".split()
['a', 'b', 'c', 'd']
>>> "a,,b".split(",")
['a', '', 'b']
>>> "".split(","), "".split()
([''], [])
>>> "a\nb\nc".splitlines()
['a', 'b', 'c']
\end{pycode}

\subsection{Combine with join()}

\code{join} is the reverse of \code{split}. It looks backward: the \emph{separator} is the object, and the
list is the argument: \code{separator.join(list)}. All items must already be strings.

%MD
%```mermaid
%flowchart LR
%    A["'a,b,c'"] -->|"split(',')"| L["['a', 'b', 'c']"]
%    L -->|"','.join(list)"| A
%```
%ENDMD
\begin{center}
\begin{tikzpicture}[font=\ttfamily\small,>=Stealth,
  ob/.style={draw=cblue,fill=cblue!8,inner sep=4pt}]
  \node[ob] (s) at (0,0) {"a,b,c"};
  \foreach \q [count=\i from 0] in {a,b,c} \node[ob,draw=cgreen,fill=cgreen!10] (l\i) at (6.2+1.3*\i,0) {"\q"};
  \node[font=\sffamily\scriptsize,text=cgray] at (7.5,0.7) {list of 3 strings};
  \draw[->,thick,cgray] (1.4,0.15) -- node[above,font=\scriptsize]{.split(",")} (5.3,0.15);
  \draw[->,thick,cgray] (5.3,-0.2) -- node[below,font=\scriptsize]{",".join(list)} (1.4,-0.2);
\end{tikzpicture}
\end{center}

\begin{pycode}
>>> ", ".join(["one", "two", "three"])
'one, two, three'
>>> "\n".join(["l1", "l2"])
'l1\nl2'
>>> ",".join([str(1), str(2)])
'1,2'
>>> ",".join([1, 2])
Traceback (most recent call last):
  ...
TypeError: sequence item 0: expected str instance, int found
\end{pycode}

\begin{gamerule}
To build one string from many pieces, collect the pieces in a list and call
\code{"".join(pieces)} once. Do not grow a string with \code{+} inside a loop.
\end{gamerule}

\subsection{Substitute with replace()}

\code{s.replace(old, new)} returns a copy with every occurrence replaced. A third argument limits the count.
It matches raw text, so it can hit the middle of words.

\begin{pycode}
>>> s = "the cat sat on the mat"
>>> s.replace("the", "a")
'a cat sat on a mat'
>>> s.replace("the", "a", 1)
'a cat sat on the mat'
>>> s.replace("at", "og")
'the cog sog on the mog'
>>> s
'the cat sat on the mat'
\end{pycode}

For whole-word or pattern replacement you need regular expressions (Chapter 17).

\subsection{Prefixes and suffixes}

\begin{pycode}
>>> f = "report_final.json"
>>> f.startswith("report"), f.endswith(".json")
(True, True)
>>> f.startswith(("a", "r"))
True
>>> f.removesuffix(".json")
'report_final'
>>> f.removeprefix("report_")
'final.json'
>>> f.removeprefix("zzz")
'report_final.json'
\end{pycode}

\code{startswith} and \code{endswith} accept a tuple of alternatives. \code{removeprefix} and \code{removesuffix}
(Python 3.9+) do nothing if the text is absent. To add a prefix or suffix, use \code{+}.

\subsection{Strip with strip()}

\code{strip()} removes characters from the \emph{ends}. With no argument it removes whitespace
(spaces, tabs, newlines). \code{lstrip()} works on the left end only, \code{rstrip()} on the right.
Given an argument, it treats it as a \emph{set of characters}, not a substring.

%MD
%```text
%  "  \n  Paris.  \n"
%   ^^^^^^^       ^^^^   <- strip() removes these
%  =>  "Paris."
%```
%ENDMD
\begin{center}
\begin{tikzpicture}[font=\ttfamily\small]
  \foreach \q [count=\i from 0] in {{\textvisiblespace},{\textvisiblespace},{\textbackslash n},{\textvisiblespace},P,a,r,i,s,.,{\textvisiblespace},{\textvisiblespace},{\textbackslash n}} {
    \ifnum\i<4 \def\fc{cred!15}\else\ifnum\i>9 \def\fc{cred!15}\else\def\fc{cgreen!15}\fi\fi
    \node[draw=cgray,fill=\fc,minimum width=7mm,minimum height=8mm] at (0.75*\i,0) {\q};
  }
  \node[font=\sffamily\scriptsize,text=cred] at (1.1,-0.8) {removed};
  \node[font=\sffamily\scriptsize,text=cred] at (7.1,-0.8) {removed};
  \node[font=\sffamily\scriptsize,text=cgreen!60!black] at (4.1,-0.8) {kept: \texttt{"Paris."}};
\end{tikzpicture}
\end{center}

\begin{pycode}
>>> raw = "  \n  Paris.  \n"
>>> raw.strip()
'Paris.'
>>> raw.lstrip()
'Paris.  \n'
>>> raw.rstrip()
'  \n  Paris.'
>>> "***hi***".strip("*")
'hi'
>>> "..ok?!".strip(".?!")
'ok'
>>> "xxabcxx".strip("xa")
'bc'
\end{pycode}

A common job is cleaning a model reply that wraps JSON in a code fence:

\begin{pycode}
>>> reply = "```json\n{\"ok\": true}\n```"
>>> reply.removeprefix("```json").removesuffix("```").strip()
'{"ok": true}'
\end{pycode}

Chaining methods works because each one returns a new string for the next to use.

\subsection{Search and select}

\begin{center}
\begin{tabularx}{\linewidth}{@{}l l X@{}}
\toprule
\textbf{Call} & \textbf{Returns} & \textbf{If not found} \\
\midrule
\code{s.find(x)} & offset (\code{int}) of the first match & \code{-1} \\
\code{s.rfind(x)} & offset of the last match & \code{-1} \\
\code{s.index(x)}, \code{s.rindex(x)} & same offsets & raises \code{ValueError} \\
\code{s.count(x)} & number of non-overlapping matches & \code{0} \\
\code{x in s} & \code{True} or \code{False} & \code{False} \\
\bottomrule
\end{tabularx}
\end{center}

\begin{pycode}
>>> t = "to be or not to be"
>>> t.find("be"), t.rfind("be"), t.find("zz")
(3, 16, -1)
>>> t.index("or")
6
>>> t.index("zz")
Traceback (most recent call last):
  ...
ValueError: substring not found
>>> t.count("to"), t.count("o")
(2, 4)
>>> t.isalnum(), "abc123".isalnum(), "12".isdigit(), "abc".isalpha()
(False, True, True, True)
\end{pycode}

\begin{trap}
\code{find} returns \code{-1} (truthy) when the text is missing and \code{0} (falsy) when it is at the start.
So \code{if s.find("x"):} is wrong. To ask ``is it there?'', use \code{"x" in s}.
\begin{pycode}
>>> "abc".find("a"), bool("abc".find("a"))
(0, False)
>>> "abc".find("z"), bool("abc".find("z"))
(-1, True)
>>> "a" in "abc", "z" in "abc"
(True, False)
\end{pycode}
\end{trap}

\subsection{Change case}

\begin{pycode}
>>> m = "hello, World"
>>> m.upper(), m.lower()
('HELLO, WORLD', 'hello, world')
>>> m.capitalize(), m.title(), m.swapcase()
('Hello, world', 'Hello, World', 'HELLO, wORLD')
>>> "Ab".lower() == "aB".lower()
True
\end{pycode}

To compare text ignoring case, lowercase both sides first.

\subsection{Set alignment}

\code{center}, \code{ljust}, \code{rjust} pad a string to a total width; an optional second argument sets the pad character.
\code{zfill} pads with zeros on the left.

\begin{pycode}
>>> "ok".center(10, "*")
'****ok****'
>>> "ok".ljust(6, ".") + "|"
'ok....|'
>>> "ok".rjust(6) + "|"
'    ok|'
>>> "7".zfill(3)
'007'
\end{pycode}

% ------------------------------------------------------------
\section{Apply Formatting}

Building text from values is called \textbf{interpolation}. Python has three styles. Know all three
so you can read old code, but write f-strings.

\begin{center}
\begin{tabularx}{\linewidth}{@{}l l X@{}}
\toprule
\textbf{Style} & \textbf{Looks like} & \textbf{Notes} \\
\midrule
Old (\code{\%}) & \code{"\%s: \%d" \% (name, n)} & C-style; values go in a tuple after \code{\%} \\
New (\code{.format}) & \code{"\{\}: \{\}".format(name, n)} & placeholders \code{\{\}}, \code{\{0\}}, \code{\{name\}} \\
f-string & \code{f"\{name\}: \{n\}"} & Python 3.6+; any expression inside braces \\
\bottomrule
\end{tabularx}
\end{center}

\subsection{Old style: \%}

The format is \code{template \% values}. A single value can go alone; several must be in a tuple.
Each \code{\%} code has a type: \code{\%s} text (works for anything), \code{\%d} integer, \code{\%f} float,
\code{\%x} hex, \code{\%e} exponent, \code{\%\%} a literal percent sign.

\begin{pycode}
>>> "%s has %d tokens" % ("prompt", 42)
'prompt has 42 tokens'
>>> "%.2f" % 3.14159
'3.14'
>>> "%5d|%-5d|" % (7, 7)
'    7|7    |'
>>> "%d%%" % 95
'95%'
>>> "%x" % 255
'ff'
\end{pycode}

\subsection{New style: format()}

\begin{pycode}
>>> "{} has {} tokens".format("prompt", 42)
'prompt has 42 tokens'
>>> "{1} before {0}".format("a", "b")
'b before a'
>>> "{name} costs {cost:.2f}".format(name="draft", cost=1.5)
'draft costs 1.50'
>>> "{0[model]}".format({"model": "demo-1"})
'demo-1'
\end{pycode}

\subsection{Newest style: f-strings}

Put \code{f} before the quote and write values or whole expressions inside braces.

\begin{pycode}
>>> model, n = "demo-1", 1234567
>>> f"{model} saw {n} tokens"
'demo-1 saw 1234567 tokens'
>>> f"{model.upper()} has {len(model)} characters"
'DEMO-1 has 6 characters'
\end{pycode}

After a colon comes a \textbf{format spec}, a tiny language shared by \code{format()} and f-strings.
Its parts, in order, are all optional:

%MD
%```text
%  f"{x:*>10,.2f}"      with x = 1234.5
%      | | | |  | |
%      | | | |  | +-- type: f float, d int, s str, x hex, b binary, % percent
%      | | | |  +---- precision: .2 = two decimals
%      | | | +------- thousands separator: , or _
%      | | +--------- width: at least 10 characters
%      | +----------- align: < left, ^ center, > right
%      +------------- fill character: *
%  result: '**1,234.50'
%```
%ENDMD
\begin{center}
\begin{tikzpicture}[font=\ttfamily\Large,>=Stealth]
  \foreach \p/\x/\col/\lab in {fill/0/cgold/{fill character (default space)},align/1/cgreen/{< left \ ^ center \ > right},width/2/cblue/{minimum width},sep/3/cpink/{, thousands separator},prec/4/cpurple/{.n decimals},type/5/cred/{f d s x b \%}} {
    \node[draw=\col,fill=\col!12,minimum width=11mm,minimum height=9mm] (n\x) at (1.5*\x,0) {};
  }
  \node at (0,0) {*}; \node at (1.5,0) {>}; \node at (3,0) {10}; \node at (4.5,0) {,}; \node at (6,0) {.2}; \node at (7.5,0) {f};
  \node[anchor=east] at (-0.8,0) {\{x:};
  \node[anchor=west] at (8.3,0) {\}};
  \node[font=\sffamily\scriptsize,text=cgold!70!black,below=3mm of n0,align=center] {fill\\char};
  \node[font=\sffamily\scriptsize,text=cgreen!70!black,below=3mm of n1,align=center] {align\\\texttt{< \string^ >}};
  \node[font=\sffamily\scriptsize,text=cblue,below=3mm of n2,align=center] {min\\width};
  \node[font=\sffamily\scriptsize,text=cpink,below=3mm of n3,align=center] {1000s\\sep};
  \node[font=\sffamily\scriptsize,text=cpurple,below=3mm of n4,align=center] {decimals};
  \node[font=\sffamily\scriptsize,text=cred,below=3mm of n5,align=center] {type};
\end{tikzpicture}
\end{center}

\begin{pycode}
>>> model, n = "demo-1", 1234567
>>> f"{1234.5:*>10,.2f}"
'**1,234.50'
>>> f"{n:,}"
'1,234,567'
>>> f"{0.4567:.1%}"
'45.7%'
>>> f"{model:>10}|{model:<10}|{model:^10}|"
'    demo-1|demo-1    |  demo-1  |'
>>> f"{7:03d}"
'007'
>>> f"{255:#x} {5:08b}"
'0xff 00000101'
\end{pycode}

\begin{center}
\begin{tabular}{@{}l l l@{}}
\toprule
\textbf{Spec} & \textbf{Meaning} & \textbf{\code{f"\{x:SPEC\}"} for \code{x = 3.14159}} \\
\midrule
\code{.2f} & 2 decimals & \code{3.14} \\
\code{8.2f} & width 8, 2 decimals & \code{'    3.14'} \\
\code{<8.2f} & left-aligned & \code{'3.14    '} \\
\code{08.2f} & zero padded & \code{00003.14} \\
\code{.2e} & scientific & \code{3.14e+00} \\
\code{.0\%} & percent & \code{314\%} \\
\bottomrule
\end{tabular}
\end{center}

Useful extras: \code{=} for debugging, \code{!r} for the \code{repr}, a nested field for a dynamic width, and
doubled braces for literal braces.

\begin{pycode}
>>> model, n = "demo-1", 1234567
>>> f"{model=}"
"model='demo-1'"
>>> f"{n * 2 = }"
'n * 2 = 2469134'
>>> f"{model!r}"
"'demo-1'"
>>> width = 8
>>> f"{model:>{width}}"
'  demo-1'
>>> f'{{"model": "{model}"}}'
'{"model": "demo-1"}'
\end{pycode}

\begin{trap}
Braces are special in f-strings. To write a literal \code{\{} or \code{\}}, double it. This matters when you
build JSON text or a prompt that contains braces. A single \code{\{} that you meant literally makes
Python try to evaluate what follows it. For real JSON, prefer \code{json.dumps(data)}.
\end{trap}

\begin{pycode}
>>> topic = "tokenizers"
>>> prompt = f"""Explain {topic}.
... Use two sentences."""
>>> prompt
'Explain tokenizers.\nUse two sentences.'
\end{pycode}

\begin{gamerule}
Use f-strings for new code. Never build prompts by gluing user text into a template with
\code{+} and no check: the text may contain braces, quotes or instructions. Keep the template
fixed and insert values deliberately.
\end{gamerule}

\section{More String Things}

Python has many more string methods (\code{partition}, \code{casefold}, \code{encode}, \code{isupper}, \ldots).
The standard documentation lists them all. Chapter 17 adds Unicode and regular expressions.

\section{Review}

A string is an immutable sequence of characters. Index it with \code{[i]}, cut it with
\code{[start:stop:step]}, and remember that methods return new strings. Use \code{split} and \code{join}
to move between text and lists, \code{in} to ask whether something is present, and f-strings to build text.

\begin{pycode}
>>> s1 = "abcdef"
>>> s1[:2] + "!" + s1[-2:]
'ab!ef'
\end{pycode}

\begin{practice}
\begin{enumerate}[itemsep=1pt]
\item What do \code{"python"[1:4]}, \code{"python"[-3:]} and \code{"python"[::-1]} return? Check by hand first.
\item Why does \code{word[0] = "T"} fail, and what are two ways to get \code{"Tokenize"} from \code{"tokenize"}?
\item A reply contains \code{"  Yes.\textbackslash n"}. Which method returns just \code{"Yes."}? What does the original value look like afterwards?
\item Why is \code{if s.find("key"):} a bug? What is the fix?
\item Turn \code{["a", "b", "c"]} into \code{"a-b-c"}, and back into a list.
\item Write an f-string that prints the float \code{0.8765} as a percentage with one decimal, right-aligned in 8 characters.
\end{enumerate}
\end{practice}

\begin{achievement}{Text Wrangler}
\begin{itemize}[itemsep=1pt]
\item Create strings with quotes, triple quotes, raw strings and f-strings, and explain echo versus print.
\item Index and slice any string and predict the result from the cut-line picture.
\item Explain why strings are immutable and keep the result of every method.
\item Clean and search text with \code{strip}, \code{split}, \code{join}, \code{replace}, \code{in} and \code{find}.
\item Format numbers and layouts with the format-spec mini-language.
\end{itemize}
\end{achievement}
